
%(BEGIN_QUESTION)
% Copyright 2006, Tony R. Kuphaldt, released under the Creative Commons Attribution License (v 1.0)
% This means you may do almost anything with this work of mine, so long as you give me proper credit

En prøve ``bordsalt'' har massen 1 kg. Hvor mange mol salt tilsvarer dette?

\vskip 10pt

Anta nå at akkurat denne mengden salt løses fullstendig i vann, slik at sluttvolumet av saltløsningen blir 850 liter. Beregn {\it molariteten} til saltløsningen.

\vskip 20pt \vbox{\hrule \hbox{\strut \vrule{} {\bf Forslag til sokratisk diskusjon} \vrule} \hrule}

\begin{itemize}
\item{} Identifiser en analyseteknologi som enkelt kan estimere molariteten i en saltløsning, og beskriv ulike varianter av denne teknologien.
\end{itemize}

\underbar{file i00568}
%(END_QUESTION)





%(BEGIN_ANSWER)

``Bordsalt'' er natriumklorid (NaCl), med 1 natriumatom bundet til 1 kloratom. Antall atommasseenheter (amu) i hvert natriumkloridmolekyl er summen av atommassene:

\vskip 10pt

\begin{itemize}
\item{} {\it Hvert molekyl NaCl inneholder:}
\item{} 1 atom Na på 22.99 amu
\item{} 1 atom Cl på 35.45 amu
\end{itemize} 

[(1 atom)(22.99 amu/atom) + (1 atom)(35.45 amu/atom)] = 58.44 g per mol NaCl

\vskip 10pt

Når vi kjenner gram per mol for NaCl, kan vi beregne hvor mange mol som trengs for å utgjøre 1 kg (1000 g) salt:

\vskip 10pt

$$\left( 1000 \hbox{ g} \over 1 \right) \left(1 \hbox{ mol NaCl} \over 58.44 \hbox{ g} \right) = 17.11 \hbox{ mol NaCl}$$

\vskip 10pt

Dermed tilsvarer 1 kg ``bordsalt'' 17.11 mol.

\vskip 10pt

{\it Molaritet}, eller {\it molarkonsentrasjon}, er antall mol løst stoff per liter løsning. Med 17.11 mol løst stoff og 850 liter løsning blir molariteten:

$${{17.11 \hbox{ mol NaCl}} \over {850 \hbox{ L}}} = 0.0201 M$$

Molaritet for et stoff representeres ofte ved å sette kjemisk symbol i hakeparentes. Selve enheten molaritet skrives som kursiv $M$. I denne oppgaven er derfor [NaCl] = 0.0201 $M$.

%(END_ANSWER)





%(BEGIN_NOTES)


%INDEX% Kjemi, støkiometri: mol
%INDEX% Kjemi: molaritet

%(END_NOTES)

